\documentclass[../DoAn.tex]{subfiles}
\begin{document}

\section{Đặt vấn đề}
\label{section:1.1}
Trong bối cảnh cuộc cách mạng công nghiệp 4.0, thị trường lao động Việt Nam đang chứng kiến sự thay đổi mạnh mẽ về cả quy mô nhu cầu tuyển dụng lẫn phương thức kết nối giữa doanh nghiệp và người lao động. Sự phát triển của các nền tảng số giúp thông tin việc làm được lan truyền nhanh hơn, tuy nhiên quá trình tuyển dụng vẫn còn tồn tại nhiều thách thức đáng kể, trong đó nổi bật là (i) nhà tuyển dụng phải xử lý số lượng lớn hồ sơ ứng viên trong thời gian ngắn, dẫn đến chi phí sàng lọc cao, độ chính xác không đồng đều và nguy cơ bỏ sót ứng viên phù hợp; (ii) ứng viên, đặc biệt là sinh viên mới ra trường hoặc người ít kinh nghiệm, thường thiếu môi trường luyện tập phỏng vấn sát với yêu cầu thực tế của từng vị trí, do đó khó nhận diện điểm mạnh, điểm yếu và cải thiện năng lực trình bày trước nhà tuyển dụng.

Những vấn đề này làm cho quá trình kết nối cung cầu lao động chưa đạt được mức hiệu quả kỳ vọng, dù dữ liệu tuyển dụng và hành vi tìm việc ngày càng được số hóa.

Các phương pháp tuyển dụng truyền thống thường phụ thuộc nhiều vào thao tác thủ công của bộ phận nhân sự, từ đọc hồ sơ, đối chiếu yêu cầu công việc, đánh giá mức độ phù hợp cho đến phản hồi cho ứng viên. Cách tiếp cận này tốn nhiều thời gian và chi phí, đồng thời khó bảo đảm tính nhất quán khi số lượng hồ sơ tăng nhanh.

Ở phía ứng viên, việc chuẩn bị phỏng vấn thường dựa trên kinh nghiệm cá nhân hoặc các nguồn tài liệu chung, chưa gắn trực tiếp với mô tả công việc, hồ sơ cá nhân và bối cảnh tuyển dụng cụ thể. Vì vậy, khoảng cách giữa yêu cầu của doanh nghiệp và khả năng tự chuẩn bị của ứng viên vẫn là một nguyên nhân quan trọng làm giảm hiệu quả tuyển dụng.

Để giải quyết các thách thức nêu trên, việc phát triển một hệ thống tích hợp công nghệ Trí tuệ nhân tạo nhằm tối ưu hóa quy trình kết nối giữa nhà tuyển dụng và ứng viên trở thành một nhu cầu cấp thiết. Giải pháp này hướng tới mục tiêu kép: một mặt, tự động hóa và nâng cao độ chính xác trong giai đoạn sàng lọc hồ sơ của nhà tuyển dụng; mặt khác, cung cấp công cụ hỗ trợ ứng viên hoàn thiện kỹ năng phỏng vấn dựa trên các kịch bản mô phỏng bám sát mô tả công việc và nội dung CV.

Nghiên cứu không chỉ góp phần tối ưu hóa hiệu suất quản trị nguồn nhân lực trong doanh nghiệp mà còn hỗ trợ nâng cao tính minh bạch, bền vững và hiệu quả vận hành của thị trường lao động.

\section{Mục tiêu và phạm vi đề tài}
\label{section:1.2}
Từ những vấn đề đã nêu, đề tài hướng tới mục tiêu xây dựng nền tảng web ``JobBridge AI'' nhằm ứng dụng Trí tuệ nhân tạo vào các hoạt động cốt lõi của quy trình tuyển dụng. Hệ thống tập trung vào hai nhóm người dùng chính là ứng viên và nhà tuyển dụng.

Đối với ứng viên, hệ thống cung cấp môi trường tìm kiếm việc làm, quản lý hồ sơ cá nhân, nộp hồ sơ ứng tuyển và luyện tập phỏng vấn với trợ lý AI dựa trên mô tả công việc đã ứng tuyển. Đối với nhà tuyển dụng, hệ thống hỗ trợ tạo và quản lý tin tuyển dụng, theo dõi danh sách ứng viên theo từng vị trí, cập nhật trạng thái ứng tuyển và sử dụng AI để đánh giá mức độ phù hợp của CV so với yêu cầu công việc.

Phạm vi hiện thực của đề tài bao gồm xây dựng ứng dụng web hoàn chỉnh với các chức năng xác thực và phân quyền theo vai trò, quản lý hồ sơ ứng viên, quản lý công ty, quản lý tin tuyển dụng, ứng tuyển, theo dõi trạng thái hồ sơ và các tính năng AI như Interview Coach, Interview Quiz và đánh giá CV cho nhà tuyển dụng.

Bên cạnh chức năng nghiệp vụ, đề tài cũng xây dựng quy trình triển khai và giám sát ở mức đủ để kiểm chứng hệ thống trong môi trường cloud. Các chức năng như thông báo thời gian thực, gợi ý việc làm bằng vector embedding hoặc huấn luyện mô hình AI riêng được xác định là hướng phát triển tiếp theo, chưa nằm trong phạm vi hoàn thiện chính của đồ án.

\section{Định hướng giải pháp}
\label{section:1.3}
Để đạt được mục tiêu trên, đề tài lựa chọn hướng xây dựng hệ thống web theo kiến trúc nhiều dịch vụ, trong đó các miền nghiệp vụ chính được tách biệt nhằm dễ phát triển, kiểm thử và mở rộng. Giao diện người dùng được tổ chức theo từng vai trò, còn dữ liệu nghiệp vụ được lưu trữ theo mô hình linh hoạt để phù hợp với đặc thù hồ sơ tuyển dụng, nội dung CV và trạng thái ứng tuyển. Các lựa chọn công nghệ cụ thể cho backend, frontend, cơ sở dữ liệu và triển khai được trình bày tập trung ở Chương 2 nhằm tránh lặp lại trong phần giới thiệu.

Đối với các tính năng AI, hệ thống không huấn luyện mô hình riêng mà khai thác mô hình ngôn ngữ lớn thông qua API để sinh phản hồi dựa trên dữ liệu nghiệp vụ sẵn có. Các ngữ cảnh như hồ sơ ứng viên, nội dung CV đã trích xuất, mô tả công việc và trạng thái ứng tuyển được sử dụng để bảo đảm phản hồi AI bám sát từng tình huống cụ thể. Đồng thời, các chức năng AI đều được đặt sau cơ chế xác thực và kiểm tra quyền để hạn chế truy cập sai phạm vi dữ liệu.

Về triển khai, đề tài định hướng xây dựng quy trình phát hành tự động và có khả năng tái tạo môi trường, thay vì triển khai thủ công từng thành phần. Phần cơ sở lý thuyết và mô tả chi tiết về container hóa, Kubernetes, GitOps, giám sát và quản lý bí mật được trình bày ở Chương 2; Chương 4 tập trung vào kết quả hiện thực và kiểm chứng trên môi trường triển khai.

\section{Bố cục đồ án}
\label{section:1.4}
Phần còn lại của báo cáo đồ án tốt nghiệp này được tổ chức như sau. Chương 2 trình bày cơ sở lý thuyết và các công nghệ nền tảng được sử dụng trong đề tài, bao gồm kiến trúc microservices, các công nghệ backend và frontend, nền tảng container hóa, Kubernetes, CI/CD GitOps và mô hình ngôn ngữ lớn. Chương 3 đi sâu vào phân tích yêu cầu và thiết kế hệ thống, bao gồm tác nhân, use case, đặc tả kịch bản, luồng nghiệp vụ, cơ sở dữ liệu và giao diện.

Chương 4 trình bày triển khai hệ thống và đánh giá thực nghiệm trên môi trường đã triển khai. Chương 5 phân tích các giải pháp kỹ thuật và đóng góp nổi bật ở góc độ hệ thống, bao gồm hạ tầng cloud có thể tái tạo bằng Terraform, quản lý bí mật bằng Azure Key Vault, tự động hóa phát hành bằng CI/CD GitOps, autoscaling bằng Kubernetes và giám sát tập trung bằng Prometheus/Grafana. Cuối cùng, Chương 6 tổng kết kết quả đạt được, đánh giá hạn chế và đề xuất hướng phát triển trong tương lai.

\end{document}
